\begin{abstract}
Spiking neural networks (SNNs) on event cameras are deployed precisely where the
standard out-of-distribution (OOD) toolkit cannot run: on neuromorphic silicon with no
floating-point logits, no gradients, and no budget for a second forward pass. We observe
that the sub-threshold membrane potential $V_{\mathrm{mem}}(t)$ of the Parametric
Leaky-Integrate-and-Fire (PLIF) neurons is a \emph{free byproduct} of the forward pass --
a register the substrate already maintains -- and that its per-channel statistics
$\varphi=[\mu,\sigma^2,\kappa]$ (a $2112$-D descriptor over $704$ channels $\times$ $4$
layers) form an OOD signal at \emph{zero additional compute}. We make four contributions.
(i)~\textbf{Gen1-C}, a reproducible, seeded event-camera corruption benchmark (six
sensor-grounded corruptions $\times$ five severities) plus a released membrane-statistics
feature set, so OOD detectors can be compared without re-running SNN inference.
(ii)~A closed-form leaky-integrator theory that predicts, per corruption, which membrane
moment moves and therefore which corruptions are detectable -- and which are
\emph{provably} not. (iii)~A \emph{motivation} experiment: feeding these corruptions to a
pretrained state-of-the-art hybrid SNN--ANN detector (no retraining) collapses its mAP,
establishing that event-stream OOD is a real failure mode worth detecting. (iv)~The
Manifold-Decomposition Detector (MDD), which decomposes each frame's deviation into
orthogonal manifold axes scored two-sided, taking the contraction corruptions above chance
for the first time. We report honest negatives -- \texttt{polarity\_flip} (a network
invariance) and \texttt{spatial\_dropout} (a signature global pooling discards) -- and
frame them as an open challenge: the Achilles' heel of event-based perception.
\advisor{framing agreed -- dataset/benchmark as a contribution; report hard corruptions as a discovered open challenge.}
\end{abstract}
